\begin{statementbox}
\textbf{统一定义的 $e=1$ 情形。} 固定点 $F$ 不在固定直线 $\ell$ 上，$H$ 为 $P$ 在 $\ell$ 上的垂足。
\[
\boxed{PF=d(P,\ell)=PH}\quad\Longleftrightarrow\quad P\text{ 在相应抛物线上}.
\]
满足条件的点不在准线上，故 $PF/PH=1$ 有意义。
这一焦准定义也是抛物线通常采用的几何定义；“第二定义”的称呼不意味着另有一个不同的等距性质。

令 $D$ 为焦点在准线上的垂足，$p=FD>0$，顶点为 $FD$ 中点 $V$。
\textbf{只有在顶点取原点、轴为 $x$ 轴且向焦点为正向时，}才有
\[
F(p/2,0),\quad \ell:x=-p/2,\quad \boxed{y^2=2px}.
\]
该标准位置下 $PF=x+p/2\ge p/2$，等号仅在顶点成立。
\end{statementbox}
